% Reference entries, W (AN-11).  Law 1: os/iword/core/edit.c
% iw_reflow_para (the space that ends a full line becomes the soft return;
% a line holds (10200 - R - L) / 120 + 1 characters, 66 at the standard
% margins: spec/iword/wp5_format.h WRAP, MINT01, MINT22; a word longer than
% the line is broken by the layout only), face_a.c (insertion: every
% character typed goes in at the cursor; Ins does nothing).
% GATES: test-pubs-iword test-iword-core test-iword-wp5-roundtrip test-iword-face test-flair-iword
\Gates{test-pubs-iword test-iword-core test-iword-wp5-roundtrip test-iword-face test-flair-iword}
\begin{Feature}{Word Wrap}
When the line you are typing is full, InitechWord moves the word that does
not fit to the start of the next line.  You keep typing and press Enter only
at the end of a paragraph.

Under the standard one inch margins a line holds 66 characters.  The space
before the word that moves down becomes a soft return (see \InRef{Return,
Soft}).

\BeginReplay{ref-wrap}
For example, type this paragraph without pressing Enter:
\Type{Initech word processing is here. The quick brown fox jumps over the lazy dog.}
The paragraph wraps after \textit{over}:
\WordScreen[rows=2]{ref-wrap}{A paragraph that wraps}
\EndReplay

\Notes
\Topic{Inserting and Deleting}
Text you type goes in at the cursor, in front of what is there; nothing is
typed over.  After every change InitechWord wraps the paragraph again, so
its lines always end where they should.

\Topic{Long Words}
A word longer than a whole line is broken where the line ends; the rest of
it continues on the next line.

\Topic{Margins}
The line length follows the margins in effect (see \InRef{Margins, Left and
Right}).  Indented lines wrap at the indent (see \InRef{Indent}).

\SeeAlso{Margins, Left and Right; Return, Hard; Return, Soft}
\end{Feature}
